\chapter{Packages and categorical structure}

This chapter develops the package layer that sits on top of the program monad
$\mathsf{SPComp}$ (Definition~\ref{def:spcomp}). Two representations coexist. A
\emph{shallow} layer stores oracle-free typed functions and inherits its
categorical structure directly from the Kleisli category of $\mathsf{SPComp}$.
A \emph{deep} layer represents computations as syntax trees (a free monad),
where oracle calls are explicit and real linking is oracle substitution. The
package-algebra results — correctness and associativity of deep linking, the
interchange law, and the reduction lemma for advantage — all live in the deep
layer.

\section{Interfaces and raw packages (shallow layer)}

\begin{definition}[Operation signature]
  \label{def:opsig}
  \lean{CatCrypt.Package.OpSig}
  \leanok
  An \emph{operation signature} $\mathsf{OpSig}$ is a pair of types
  $(\mathsf{src}, \mathsf{tgt})$: the input and output type of a single oracle
  operation.
\end{definition}

\begin{definition}[Interface]
  \label{def:interface}
  \uses{def:opsig}
  \lean{CatCrypt.Package.Interface}
  \leanok
  An \emph{interface} is a finite set of entries, each pairing an operation
  identifier $\mathrm{id} \in \mathbb{N}$ with a signature, subject to
  uniqueness of identifiers. Interfaces carry an empty interface, a membership
  predicate $\mathsf{has}\,I\,\mathrm{id}\,S\,T$, a disjointness relation, and a
  union of identifier-disjoint interfaces.
\end{definition}

\begin{definition}[Typed function]
  \label{def:typedfun}
  \uses{def:spcomp}
  \lean{CatCrypt.Package.TypedFun}
  \leanok
  A \emph{typed function} bundles an input type $\mathsf{src}$, an output type
  $\mathsf{tgt}$, and an implementation $\mathsf{impl} : \mathsf{src} \to
  \mathsf{SPComp}\,\mathsf{tgt}$. These implementations are oracle-free: they
  are ordinary stateful-probabilistic computations with nothing left to
  substitute.
\end{definition}

\begin{definition}[Raw package]
  \label{def:rawpackage}
  \uses{def:typedfun}
  \lean{CatCrypt.Package.RawPackage}
  \leanok
  A \emph{raw package} is a finite set of package entries, each pairing an
  operation identifier with a typed function, again with unique identifiers. It
  carries $\mathsf{ids}$, $\mathsf{has}$, $\mathsf{lookup}$, a $\mathsf{singleton}$
  constructor, and an empty package.
\end{definition}

\begin{definition}[Parallel composition of raw packages]
  \label{def:rawpackage-par}
  \uses{def:rawpackage}
  \lean{CatCrypt.Package.RawPackage.par}
  \leanok
  For identifier-disjoint packages $P_1, P_2$, the parallel composition
  $\mathsf{par}\,P_1\,P_2$ is the union of their entries; disjointness of
  identifiers guarantees the uniqueness invariant.
\end{definition}

\begin{definition}[Resolving an operation]
  \label{def:rawpackage-resolve}
  \uses{def:rawpackage, def:spcomp}
  \lean{CatCrypt.Package.RawPackage.resolve}
  \leanok
  $\mathsf{resolve}\,P\,\mathrm{id}\,x$ looks up the implementation of
  operation $\mathrm{id}$, checks that the stored source and target types match
  the requested ones, and calls it; if the operation is absent or the types
  mismatch it returns $\mathsf{SPComp.fail}$. This is how a shallow package is
  run.
\end{definition}

\begin{definition}[Linking of raw packages is a stub]
  \label{def:rawpackage-link}
  \uses{def:rawpackage}
  \lean{CatCrypt.Package.RawPackage.link}
  \leanok
  $\mathsf{link}\,P_1\,P_2 = P_1$ is a \emph{shallow-layer no-op stub}. Because
  $\mathsf{RawPackage}$ stores oracle-free typed functions, there are no oracle
  calls to substitute, so linking has nothing to do. Oracle-substituting
  linking is defined in the deep embedding
  (Definition~\ref{def:deeppackage-link}) and proved correct and associative by
  Theorems~\ref{thm:runpkg-link} and~\ref{thm:runpkg-link-assoc}.
\end{definition}

\section{The Kleisli category of \texorpdfstring{$\mathsf{SPComp}$}{SPComp}}

\begin{definition}[Kleisli category of $\mathsf{SPComp}$]
  \label{def:klspcomp}
  \uses{def:spcomp, def:spcomp-ops}
  \lean{CatCrypt.Core.KlSPComp}
  \leanok
  $\mathsf{KlSPComp}$ is the type of objects of the Kleisli category of
  $\mathsf{SPComp}$ (an opaque copy of $\mathsf{Type}$). A morphism $\alpha \to
  \beta$ is a function $\alpha \to \mathsf{SPComp}\,\beta$; identity is
  $\mathsf{pure}$ and composition is $\mathsf{bind}$. The category laws are the
  monad laws (Lemmas~\ref{lem:pure-bind}, \ref{lem:bind-pure},
  \ref{lem:bind-assoc}). With $\mathrm{Sum}$ as tensor and $\mathrm{Empty}$ as
  unit, $\mathsf{KlSPComp}$ is a symmetric monoidal category.
\end{definition}

\begin{lemma}[Composition is bind]
  \label{lem:klspcomp-comp}
  \uses{def:klspcomp}
  \lean{CatCrypt.Core.KlSPComp.comp_apply}
  \leanok
  For $f : \alpha \to \beta$ and $g : \beta \to \gamma$ in $\mathsf{KlSPComp}$,
  $(f \fatsemi g)\,x = \mathsf{bind}\,(f\,x)\,g$.
\end{lemma}

\begin{definition}[Cocartesian monoidal category]
  \label{def:cocartesian-class}
  \lean{CategoryTheory.CocartesianMonoidalCategory}
  \leanok
  A monoidal category is \emph{cocartesian} when its unit is initial and its
  tensor is the categorical binary coproduct (the dual of a cartesian monoidal
  category). It provides coprojections $\mathsf{inl}, \mathsf{inr}$ and a
  copairing $\mathsf{desc}$ with the expected universal property.
\end{definition}

\begin{definition}[Copairing in $\mathsf{KlSPComp}$]
  \label{def:klspcomp-codesc}
  \uses{def:klspcomp}
  \lean{CatCrypt.Core.KlSPComp.klDesc}
  \leanok
  Given $f : \alpha \to \gamma$ and $g : \beta \to \gamma$ in
  $\mathsf{KlSPComp}$, the copairing $\mathsf{klDesc}\,f\,g$ maps
  $\mathrm{inl}\,a \mapsto f\,a$ and $\mathrm{inr}\,b \mapsto g\,b$. The
  coprojections $\mathsf{klInl}, \mathsf{klInr}$ embed a summand via
  $\mathsf{pure}$.
\end{definition}

\begin{proposition}[$\mathsf{KlSPComp}$ is cocartesian]
  \label{prop:klspcomp-cocartesian}
  \uses{def:cocartesian-class, def:klspcomp-codesc}
  \lean{CatCrypt.Core.KlSPComp.binaryCoproductIsColimit}
  \leanok
  The cofan $\mathrm{BinaryCofan}(\mathsf{klInl}, \mathsf{klInr})$ is a colimit:
  $\mathrm{Sum}$ is the coproduct in $\mathsf{KlSPComp}$, with $\mathrm{Empty}$
  initial. Hence $\mathsf{KlSPComp}$ is a cocartesian monoidal category.
\end{proposition}

\begin{definition}[Affine monoidal category]
  \label{def:affine-class}
  \lean{CategoryTheory.AffineMonoidalCategory}
  \leanok
  An \emph{affine} monoidal category equips every object with a natural discard
  morphism $\mathsf{del}\,X : X \to \mathbf{1}$ satisfying $\mathsf{del}\,
  \mathbf{1} = \mathrm{id}$ and naturality $f \fatsemi \mathsf{del}\,Y =
  \mathsf{del}\,X$. Every computation can be discarded.
\end{definition}

\begin{definition}[Discard in $\mathsf{KlSPComp}$]
  \label{def:klspcomp-del}
  \uses{def:affine-class, def:klspcomp}
  \lean{CatCrypt.Core.KlSPComp.klDel}
  \leanok
  The discard morphism $\mathsf{klDel}\,\alpha : \alpha \to \mathrm{Empty}$
  sends every value to the failing computation $\mathsf{SPComp.fail}$. This
  makes $\mathsf{KlSPComp}$ an affine monoidal category.
\end{definition}

\section{The family (bi)category and package interfaces}

\begin{definition}[Category of families]
  \label{def:famobj}
  \lean{CatCrypt.Category.FamObj}
  \leanok
  For a category $C$, an object of $\mathsf{FamObj}\,C$ is a pair $(\iota,
  \mathsf{objs})$ of an index type $\iota$ and a family $\mathsf{objs} : \iota
  \to C$. A morphism carries an index equivalence $\sigma : \iota_1 \simeq
  \iota_2$ together with componentwise $C$-morphisms $\mathsf{map}\,i :
  \mathsf{objs}_1\,i \to \mathsf{objs}_2(\sigma\,i)$. This is a symmetric
  monoidal category with $\mathrm{Sum}$ tensor and $\mathrm{Empty}$ unit.
\end{definition}

\begin{definition}[Package interface]
  \label{def:pkginterface}
  \uses{def:klspcomp}
  \lean{CatCrypt.Category.PkgInterface}
  \leanok
  A \emph{package interface} $I$ is an index type $\iota$ together with domain
  and codomain families $\mathsf{doms}, \mathsf{codoms} : \iota \to
  \mathsf{KlSPComp}$. It presents domain and codomain objects
  $\mathsf{domFam}\,I, \mathsf{codomFam}\,I$ in $\mathsf{FamObj}\,
  \mathsf{KlSPComp}$. Parallel combination $\mathsf{tensor}$ (via $\mathrm{Sum}$)
  and the empty interface $\mathsf{unit}$ agree definitionally with the family
  tensor and unit.
\end{definition}

\begin{definition}[Package implementation]
  \label{def:pkgimpl}
  \uses{def:pkginterface, def:famobj}
  \lean{CatCrypt.Category.PkgImpl}
  \leanok
  A \emph{package implementation} of interface $I$ is a morphism
  $\mathsf{domFam}\,I \to \mathsf{codomFam}\,I$ in
  $\mathsf{FamObj}\,\mathsf{KlSPComp}$: an index permutation together with, for
  each operation $i$, a Kleisli arrow $\mathsf{doms}\,i \to
  \mathsf{SPComp}(\mathsf{codoms}(\sigma\,i))$.
\end{definition}

\begin{definition}[Kleisli linking of package implementations]
  \label{def:pkgimpl-link}
  \uses{def:pkgimpl}
  \lean{CatCrypt.Category.PkgImpl.link}
  \leanok
  Linking $\mathsf{link}\,p_1\,p_2 = p_1 \fatsemi p_2$ is composition in
  $\mathsf{FamObj}\,\mathsf{KlSPComp}$. At each operation it computes as
  $\mathsf{bind}\,(p_1.\mathsf{map}\,i\,x)\,(p_2.\mathsf{map}(p_1.\sigma\,i))$:
  a first-order \emph{data pipeline} (run $p_1$, feed the result to $p_2$),
  \emph{not} second-order oracle substitution. Cryptographic package linking is
  the oracle-resolving $\mathsf{DeepPackage.link}$
  (Definition~\ref{def:deeppackage-link}).
\end{definition}

\begin{definition}[Parallel composition of package implementations]
  \label{def:pkgimpl-par}
  \uses{def:pkgimpl, def:famobj}
  \lean{CatCrypt.Category.PkgImpl.par}
  \leanok
  $\mathsf{par}\,p_1\,p_2 = p_1 \otimes p_2$ is the family tensor of the two
  implementations. It dispatches an operation $\mathrm{inl}\,i$ to $p_1$ and
  $\mathrm{inr}\,j$ to $p_2$. This is parallel composition proper, shared with
  the package category, and inherits the symmetric monoidal laws from
  $\mathsf{FamObj}$.
\end{definition}

\begin{definition}[Bridge from deep interfaces to package interfaces]
  \label{def:deepinterface-topkg}
  \uses{def:pkginterface, def:deepinterface}
  \lean{CatCrypt.Category.DeepInterface.toPkgInterface}
  \leanok
  A list-based $\mathsf{DeepInterface}$ becomes a family-based
  $\mathsf{PkgInterface}$ by $\mathrm{Fin}$-indexing its operation list: each
  position becomes an index, with domain and codomain read off the entry.
\end{definition}

\section{The deep embedding}

\begin{definition}[Raw code: a free monad]
  \label{def:rawcode}
  \uses{def:spcomp}
  \lean{CatCrypt.Deep.RawCode.inductionOn}
  \leanok
  $\mathsf{RawCode}\,\alpha$ is the free monad of computation syntax trees,
  with constructors $\mathsf{ret}$, $\mathsf{bind}$, $\mathsf{sample}$,
  $\mathsf{get}$, $\mathsf{put}$, $\mathsf{fail}$, $\mathsf{oracleCall}$, and
  $\mathsf{embed}$ (which injects an arbitrary $\mathsf{SPComp}$ value as opaque
  code, expressing completeness). Unlike the shallow layer, oracle calls are
  first-class syntax. The structural recursor $\mathsf{inductionOn}$ enumerates
  these constructors.
\end{definition}

\begin{definition}[Deep interface]
  \label{def:deepinterface}
  \lean{CatCrypt.Deep.DeepInterface}
  \leanok
  A \emph{deep interface} is a list of operation entries $(\mathrm{id}, \mathrm{dom},
  \mathrm{codom})$ of matching universe. It carries an empty interface, a
  membership predicate $\mathsf{contains}$, and extensionality by the operation
  list.
\end{definition}

\begin{definition}[Valid code bundle]
  \label{def:validcodebundle}
  \uses{def:rawcode}
  \lean{CatCrypt.Deep.ValidCodeBundle}
  \leanok
  Validity of code is the inductive predicate $\mathsf{ValidCode}\,L$ that every
  $\mathsf{get}/\mathsf{put}$ touches only a location id in the declared set
  $L$. A \emph{valid code bundle} pairs a $\mathsf{RawCode}\,\alpha$ with its
  $\mathsf{ValidCode}\,L$ proof; by proof irrelevance two bundles are equal
  whenever their code is.
\end{definition}

\begin{definition}[Deep package]
  \label{def:deeppackage}
  \uses{def:deepinterface, def:validcodebundle}
  \lean{CatCrypt.Deep.DeepPackage}
  \leanok
  A \emph{deep package} bundles a finite location set $\mathsf{locs}$, an
  import interface, an export interface, and, for every exported operation, an
  implementation mapping a domain value to a valid code bundle over
  $\mathsf{locs}$. The location invariant makes modular composition and
  separation reasoning sound.
\end{definition}

\begin{definition}[Separation]
  \label{def:deeppackage-sep}
  \uses{def:deeppackage}
  \lean{CatCrypt.Deep.DeepPackage.sep}
  \leanok
  Two deep packages are \emph{separated}, $\mathsf{sep}\,p_1\,p_2$, when their
  location sets are disjoint. Separated packages cannot interfere with each
  other's state.
\end{definition}

\begin{definition}[Oracle substitution]
  \label{def:substoracle}
  \uses{def:rawcode}
  \lean{CatCrypt.Deep.RawCode.substOracle}
  \leanok
  $\mathsf{substOracle}\,c\,\mathsf{env}$ recursively replaces every
  $\mathsf{oracleCall}\,\mathrm{op}\,\mathrm{dom}\,\mathrm{codom}\,x$ node in
  $c$ by $\mathsf{env}\,\mathrm{op}\,\mathrm{dom}\,\mathrm{codom}\,x$, leaving
  the other constructors intact. This is the engine of real package linking.
\end{definition}

\begin{lemma}[Substitution composes]
  \label{lem:substoracle-comp}
  \uses{def:substoracle}
  \lean{CatCrypt.Deep.RawCode.substOracle_comp}
  \leanok
  $(\mathsf{substOracle}\,c\,\mathsf{env}_1)\,\mathsf{substOracle}\,\mathsf{env}_2
  = \mathsf{substOracle}\,c\,(\lambda\,\mathrm{op}\,\mathrm{dom}\,\mathrm{codom}\,x.\,
  \mathsf{substOracle}\,(\mathsf{env}_1\,\mathrm{op}\,\mathrm{dom}\,\mathrm{codom}\,x)\,
  \mathsf{env}_2)$. Double substitution is a single substitution through a
  composed environment.
\end{lemma}

\begin{lemma}[Substitution is a no-op on oracle-free code]
  \label{lem:substoracle-self}
  \uses{def:substoracle}
  \lean{CatCrypt.Deep.RawCode.substOracle_eq_self}
  \leanok
  If $c$ contains no $\mathsf{oracleCall}$ node (predicate $\mathsf{NoOracleCall}\,c$),
  then $\mathsf{substOracle}\,c\,\mathsf{env} = c$ for every environment. This is
  the foundational lemma for reduction packages, whose non-delegated operations
  never call imports.
\end{lemma}

\begin{definition}[Deep linking]
  \label{def:deeppackage-link}
  \uses{def:deeppackage, def:substoracle}
  \lean{CatCrypt.Deep.DeepPackage.link}
  \leanok
  $\mathsf{link}\,p_1\,p_2$ connects $p_2$'s exports to $p_1$'s imports. It uses
  the union of locations, imports $p_2$'s imports, exports $p_1$'s exports, and
  implements each exported operation by $\mathsf{substOracle}$ on $p_1$'s code,
  replacing every call to an operation in $p_2$'s exports with $p_2$'s
  implementation (and every other call with $\mathsf{fail}$). This is the real,
  oracle-resolving package linking.
\end{definition}

\begin{definition}[Deep parallel composition]
  \label{def:deeppackage-par}
  \uses{def:deeppackage, def:deeppackage-sep}
  \lean{CatCrypt.Deep.DeepPackage.par}
  \leanok
  For separated $p_1, p_2$, $\mathsf{par}\,p_1\,p_2$ takes the union of
  locations and the concatenation of import/export interfaces, dispatching each
  exported operation to whichever package declares it. Separation is preserved
  by parallel composition on either side.
\end{definition}

\begin{definition}[Identity package]
  \label{def:deeppackage-id}
  \uses{def:deeppackage}
  \lean{CatCrypt.Deep.DeepPackage.id}
  \leanok
  $\mathsf{id}\,I$ is the stateless passthrough for interface $I$: it exports
  and imports $I$, and implements each operation by a single
  $\mathsf{oracleCall}$ forwarding the request. It is the unit of linking at the
  evaluation level.
\end{definition}

\section{Deep-to-shallow bridge and package advantage}

\begin{definition}[Evaluation to $\mathsf{SPComp}$]
  \label{def:eval}
  \uses{def:rawcode, def:spcomp}
  \lean{CatCrypt.Deep.RawCode.eval}
  \leanok
  $\mathsf{RawCode.eval}$ interprets a syntax tree as an $\mathsf{SPComp}$
  program, mapping each constructor to the corresponding monad operation and
  sending unresolved $\mathsf{oracleCall}$ nodes to $\mathsf{fail}$; a variant
  $\mathsf{evalWith}$ takes an explicit oracle handler. In the reverse
  direction, an elaboration-time reifier $\mathsf{rawCode\%}$ turns a shallow
  $\mathsf{SPComp}$ program over the six core constructors back into
  $\mathsf{RawCode}$, and the two round-trip up to the monad laws.
\end{definition}

\begin{lemma}[Substitution commutes with evaluation]
  \label{lem:eval-substoracle}
  \uses{def:substoracle, def:eval}
  \lean{CatCrypt.Deep.eval_substOracle}
  \leanok
  Evaluating $\mathsf{substOracle}\,c\,\mathsf{env}$ equals evaluating $c$ with
  the handler that evaluates each environment entry:
  $(\mathsf{substOracle}\,c\,\mathsf{env}).\mathsf{eval} =
  c.\mathsf{evalWith}(\lambda\ldots.\,(\mathsf{env}\ldots).\mathsf{eval})$. This
  is the bridge that makes syntactic linking semantically correct.
\end{lemma}

\begin{definition}[Running a deep package]
  \label{def:runpkg}
  \uses{def:deeppackage, def:eval}
  \lean{CatCrypt.Crypto.runPkg}
  \leanok
  $\mathsf{runPkg}\,p$ evaluates the conventional main export
  $(0, \mathrm{Unit}, \mathrm{Bool})$ of a deep package to an
  $\mathsf{SPComp}\,\mathrm{Bool}$ (returning $\mathsf{pure}\,\mathsf{true}$ if
  no such export exists), mirroring the cryptographic convention that an
  adversary's main procedure takes no input and returns a boolean guess.
\end{definition}

\begin{definition}[Deep nominal advantage]
  \label{def:deepnomadv}
  \uses{def:runpkg, def:deeppackage-link}
  \lean{CatCrypt.Crypto.DeepNomAdvantage}
  \leanok
  For nominal packages $G, G', A$,
  \[
    \mathsf{DeepNomAdvantage}\,G\,G'\,A =
    \mathsf{Advantage}\bigl(\mathsf{runPkg}(\mathsf{link}\,A\,G),\
    \mathsf{runPkg}(\mathsf{link}\,A\,G')\bigr),
  \]
  the distinguishing probability of adversary $A$ linked against each game via
  real (oracle-resolving) deep linking.
\end{definition}

\begin{theorem}[Correctness of linking under evaluation]
  \label{thm:runpkg-link}
  \uses{def:runpkg, def:deeppackage-link, lem:eval-substoracle}
  \lean{CatCrypt.Crypto.runPkg_link}
  \leanok
  $\mathsf{runPkg}(\mathsf{link}\,p_1\,p_2)$ equals evaluating $p_1$'s main
  procedure with the oracle handler that resolves each call through $p_2$'s
  exports (and sends absent operations to $\mathsf{fail}$). Syntactic linking on
  the free monad implements semantic oracle resolution on $\mathsf{SPComp}$.
\end{theorem}

\begin{theorem}[Associativity of linking under evaluation]
  \label{thm:runpkg-link-assoc}
  \uses{def:runpkg, def:deeppackage-link, lem:substoracle-comp, lem:eval-substoracle}
  \lean{CatCrypt.Crypto.runPkg_link_assoc}
  \leanok
  $\mathsf{runPkg}(\mathsf{link}(\mathsf{link}\,p_1\,p_2)\,p_3) =
  \mathsf{runPkg}(\mathsf{link}\,p_1\,(\mathsf{link}\,p_2\,p_3))$. The two
  packages are not definitionally equal — their oracle environments differ
  structurally — but they agree after evaluation, via $\mathsf{substOracle\_comp}$
  and the fact that $\mathsf{fail}$ substitutes to $\mathsf{fail}$.
\end{theorem}

\begin{theorem}[Interchange law]
  \label{thm:runpkg-interchange}
  \uses{def:runpkg, def:deeppackage-link, def:deeppackage-par}
  \lean{CatCrypt.Crypto.runPkg_interchange}
  \leanok
  Under an interface-matching condition on the oracle environments,
  \[
    \mathsf{runPkg}\bigl(\mathsf{link}\,(\mathsf{par}\,p_1\,p_2)\,(\mathsf{par}\,p_3\,p_4)\bigr)
    = \mathsf{runPkg}\bigl(\mathsf{par}\,(\mathsf{link}\,p_1\,p_3)\,(\mathsf{link}\,p_2\,p_4)\bigr).
  \]
  Linking after tensoring agrees with tensoring after linking, the exchange law
  that makes packages a (bi)categorical composition structure.
\end{theorem}

\begin{theorem}[Reduction lemma for advantage]
  \label{thm:deepnomadv-link}
  \uses{def:deepnomadv, thm:runpkg-link-assoc}
  \lean{CatCrypt.Crypto.DeepNomAdvantage_link}
  \leanok
  For compatible registries,
  \[
    \mathsf{DeepNomAdvantage}\,G_0\,G_1\,(\mathsf{link}\,A\,R) =
    \mathsf{DeepNomAdvantage}\,(\mathsf{link}\,R\,G_0)\,(\mathsf{link}\,R\,G_1)\,A.
  \]
  A reduction wrapper $R$ can be absorbed into the adversary — the standard move
  in hybrid arguments — and this follows directly from
  Theorem~\ref{thm:runpkg-link-assoc}.
\end{theorem}

\begin{definition}[Game-based security for deep packages]
  \label{def:nompkgsecure}
  \uses{def:deepnomadv}
  \lean{CatCrypt.Crypto.NomPkgSecure}
  \leanok
  $\mathsf{NomPkgSecure}\,G\,G'\,\varepsilon$ holds when for every adversary
  package $A$, $\mathsf{DeepNomAdvantage}\,G\,G'\,A \le \varepsilon\,A$. The
  bound is adversary-dependent, supporting concrete security; it is reflexive
  at $0$, symmetric, and transitive with pointwise-added bounds via the triangle
  inequality.
\end{definition}
